\chapter{Angular cutoffs and smooth limits}
\label{mht:ch:angular-cutoffs}

The finite reductions of Chapter~\ref{mht:ch:finite-reductions}
preserve the cyclic pairing. To remove the cutoff, their projections
must also approximate smooth fields with controlled norms. On the
round sphere the polar pairing determines such a family. Its rotation
weights pair to one, so a symmetric harmonic cutoff omits the partners
of its highest modes. Retaining one additional band restores the
pairing and gives uniform estimates for each fixed higher operation.

\section{Rotation weights and the Green polarization}

Keep the round-sphere representative
\[
 f_0=\frac{z}{1+|z|^2},\qquad
 \eta_0=d_Ef_0=-\frac{z^2}{(1+|z|^2)^2}d\bar z
\]
and the circular polar trace of
Chapter~\ref{mht:ch:polar-neumann}. The scalar cyclic projection is
the spherical mean:
\begin{equation}\label{mht:eq:angular-mean}
 p_0(f)=\int_{S^2}f\,d\mu,
 \qquad d\mu=\frac{dx\,dy}{\pi(1+|z|^2)^2},
 \qquad \mu(S^2)=1.
\end{equation}
Indeed, substitute $\eta_0$ into
$p_0(f)=\kappa^{-1}\tau_E(f\eta_0)$ and use
$dz\wedge d\bar z=-2i\,dx\wedge dy$.

As before, $V=\ker p_0\subset E^0$, $W=\ker p_1\subset E^1$,
and $d_E:V\to W$ has inverse $h$. Write
\[
 \omega(v,w)=\tau_E(vd_Ew).
\]

A function has rotation weight $m\in\mathbb Z$ if
$f(e^{i\theta}z)=e^{im\theta}f(z)$. Let $\mathcal P_m$ be the
corresponding Fourier projector and put $V_m=\mathcal P_mV$.
The polar integral gives the selection rule
\begin{equation}\label{mht:eq:angular-weight-pairing}
 \omega(V_m,V_n)=0\qquad\text{unless }m+n=1.
\end{equation}
To see the shift, differentiation by $\bar z$ raises the coefficient
weight by one. The coefficient of $vd_Ew$ therefore has weight
$m+n+1$, whereas the factor $z^{-2}$ requires weight two.

Let $V_-$ and $V_+$ be the closed subspaces with nonpositive and
positive weights, respectively. The mean condition affects weight
zero, and the first-jet condition affects weight one. Hence
\[
 V=V_-\oplus V_+.
\]
Both summands are isotropic. The Green operator identifies their
pairing with the ordinary Hermitian pairing on one summand.

\begin{proposition}\label{mht:prop:angular-polarization}
The antilinear maps
\begin{equation}\label{mht:eq:angular-green-pair}
 \mathcal Cv=h(\eta_0\overline v)\quad(v\in V_-),
 \qquad
 \mathcal Rf=\overline{\frac{d_Ef}{\eta_0}}\quad(f\in V_+)
\end{equation}
are inverse continuous isomorphisms. For $u,v\in V_-$,
\begin{equation}\label{mht:eq:angular-positive-pairing}
 \omega(u,\mathcal Cv)
       =\kappa\int_{S^2}u\overline v\,d\mu.
\end{equation}
For every integer $s\geq2$, the round Sobolev norms satisfy
\[
 \|\mathcal Rf\|_{H^s}\leq C_s\|f\|_{H^{s+6}},
 \qquad
 \|\mathcal Cv\|_{H^{s+1}}\leq C_s\|v\|_{H^s}.
\]
The loss of six derivatives is not asserted to be optimal.
\end{proposition}

\begin{proof}
The mean of $\overline v$ is zero. In the formula for
$h(\eta_0\overline v)$, the other scalar correction is the mean of
$f_0\overline v$. Its weights are strictly positive, so it also
vanishes. Thus
\[
 \mathcal Cv=G(\eta_0\overline v),\qquad
 d_E\mathcal Cv=\eta_0\overline v.
\]

This function has positive weights. Its first jet is zero by polar
Stokes and the zero mean of $v$. Equation~\eqref{mht:eq:angular-mean}
now proves \eqref{mht:eq:angular-positive-pairing}.

The quotient in \eqref{mht:eq:angular-green-pair} requires care only
at $o$. If $f\in V_+$, the coefficient of $d_Ef$ has weights at
least two. Every monomial $z^a\bar z^b$ in its Taylor polynomial
satisfies $a-b\geq2$, and hence is divisible by $z^2$.

Fix a nonnegative integer $r$. Expand that coefficient through
degree $r+2$. The remainder and each derivative of order $j\leq r$
are bounded by a constant times $|z|^{r+3-j}$. After division by
$z^2$, the derivatives of order at most $r$ tend to zero at $o$.
The divided polynomial supplies their limits. This proves smooth
division and bounds the quotient in $C^r$ by the original coefficient
in $C^{r+3}$.

In the chart $w=1/z$, one has
$\eta_0=(1+|w|^2)^{-2}d\bar w$, so division there is ordinary
smooth division. Apply the local estimate with $r=s$, differentiate
$f$ once, and use the Sobolev embedding
$H^{s+6}\hookrightarrow C^{s+4}$ in real dimension two.
This gives the estimate for $\mathcal R$. The estimate for
$\mathcal C$ follows from the elliptic gain of one derivative for $G$.

The quotient defining $\mathcal Rf$ has nonpositive weights after
conjugation. Its mean is zero because
$\tau_E(d_Ef)=\kappa j(f)=0$. Thus $\mathcal Rf\in V_-$.
The differential formula proves $\mathcal R\mathcal C=1$.
Conversely, $Gd_Ef=f-f(o)=f$ on positive weights, proving
$\mathcal C\mathcal R=1$.
\end{proof}

The normalized form $\kappa^{-1}\omega(u,\mathcal Cv)$ is positive
Hermitian on $V_-$. This statement uses the two isotropic sectors.
The original bilinear form $\omega$ remains skew.

\section{The additional harmonic band}

Use the sphere of radius one, whose scalar Laplacian has eigenvalues
$\ell(\ell+1)$. For an integer $L\geq0$, let $\mathscr H_{\leq L}$
be its harmonics of degrees at most $L$, and set
\[
 W_L=V_-\cap\mathscr H_{\leq L},\qquad
 d_L=\dim W_L=\frac{(L+1)(L+2)}2-1.
\]
There are $\ell+1$ nonpositive weights at degree $\ell$, and the
constant has been removed. In particular, $W_0=0$.

\begin{proposition}\label{mht:prop:angular-radical}
For $L\geq1$,
\begin{equation}\label{mht:eq:angular-ordinary-cutoff}
 V\cap\mathscr H_{\leq L}=W_L\oplus\mathcal CW_{L-1}.
\end{equation}
The restricted form has radical
$W_L\cap W_{L-1}^{\perp_{L^2}}$, of dimension $L+1$, and rank
$L(L+1)-2$. The subspace
\begin{equation}\label{mht:eq:angular-staggered-space}
 \widehat V_L=W_L\oplus\mathcal CW_L,
 \qquad \dim\widehat V_L=L(L+3),
\end{equation}
is nondegenerate, and the spaces $\widehat V_L$ are nested.
\end{proposition}

\begin{proof}
The form $\eta_0=d_Ef_0$ belongs to the degree-one angular
representation. Multiplication by it takes degrees at most $L-1$
to degrees at most $L$. The round Dolbeault inverse preserves each
nonconstant angular representation. Its normalization adds no
constant on positive weights. Therefore
\[
 \mathcal CW_{L-1}\subset V_+\cap\mathscr H_{\leq L}.
\]

The space on the right has dimension $L(L+1)/2-1$: the first-jet
condition removes one dimension from the positive-weight harmonics.
This equals $d_{L-1}$, so the inclusion is an equality.

Equation~\eqref{mht:eq:angular-positive-pairing} identifies the
cross pairing with the Hermitian pairing between $W_L$ and $W_{L-1}$.
Its radical is the orthogonal complement of $W_{L-1}$ in $W_L$.
The dimension difference is $d_L-d_{L-1}=L+1$, and the rank is
$2d_{L-1}=L(L+1)-2$. Retaining $\mathcal CW_L$ instead makes the
cross pairing nondegenerate. Nesting follows from $W_L\subset W_{L+1}$.
\end{proof}

At the first cutoff, take
\begin{equation}\label{mht:eq:angular-first-modes}
 v=\frac{|z|^2}{1+|z|^2}-\frac12,
 \qquad u=\frac{\bar z}{1+|z|^2}.
\end{equation}
Their partners are
\[
 \mathcal Cv=\frac{z^2\bar z}{2(1+|z|^2)^2},
 \qquad \mathcal Cu=\frac{z^2}{2(1+|z|^2)^2}.
\]

Differentiation proves the defining equations, and each partner
vanishes at $o$. All four functions are smooth on the sphere and
satisfy the Neumann condition. In the basis
$(v,u,\mathcal Cv,\mathcal Cu)$ of $\widehat V_1$, the pairing is
\begin{equation}\label{mht:eq:angular-first-pairing}
 \kappa\begin{pmatrix}
 0&0&1/12&0\\
 0&0&0&1/6\\
 -1/12&0&0&0\\
 0&-1/6&0&0
 \end{pmatrix}.
\end{equation}
The two nonzero entries are the spherical norms of $v$ and $u$.
The mode $w$ of \eqref{mht:eq:finite-two-modes} is $2\mathcal Cv+u$.
Thus the earlier six-dimensional complex retains a mixed-weight
subspace of this first ten-dimensional coefficient complex.

\section{Uniform estimates for the projections}

Let $S_L^-$ be the orthogonal spectral projector onto $W_L$.
For $f=f_-+f_+\in V$, define
\begin{equation}\label{mht:eq:angular-projection}
 \widehat\Pi_Lf=S_L^-f_-+\mathcal CS_L^-\mathcal Rf_+.
\end{equation}
The two antilinear maps make the second term complex linear.
The cross pairing \eqref{mht:eq:angular-positive-pairing} shows that
this is the $\omega$-orthogonal projection onto $\widehat V_L$.
Consequently
\[
 \widehat\Pi_L\widehat\Pi_M=\widehat\Pi_{\min(L,M)},
 \qquad
 \omega(\widehat\Pi_Lf,g)=\omega(f,\widehat\Pi_Lg).
\]

\begin{theorem}\label{mht:thm:angular-tame}
For integers $s\geq2$ and $r\geq0$, there are constants independent
of $L$ such that
\begin{align}
 \|\widehat\Pi_Lf\|_{H^s}
       &\leq C_s\|f\|_{H^{s+6}},\notag\\
 \|(1-\widehat\Pi_L)f\|_{H^s}
       &\leq C_{s,r}(1+L)^{-r}\|f\|_{H^{s+r+6}}.
 \label{mht:eq:angular-tame}
\end{align}
The projections converge to the identity uniformly on bounded
smooth sets in every smooth seminorm.
\end{theorem}

\begin{proof}
The projections to positive and nonpositive rotation weights commute
with the round Laplacian and are bounded on every Sobolev space.
Spectral expansion on $V_-$ gives
\[
 \|(1-S_L^-)v\|_{H^s}
       \leq C_r(1+L)^{-r}\|v\|_{H^{s+r}}.
\]
Indeed, every omitted harmonic has degree greater than $L$.
Apply this inequality to $f_-$ and $\mathcal Rf_+$, and use both
estimates in Proposition~\ref{mht:prop:angular-polarization}.
The same argument without the tail factor bounds the projection.
Every bounded smooth set has the higher Sobolev bounds on the right.
\end{proof}

The norms
\begin{equation}\label{mht:eq:angular-adapted-norm}
 \|f\|_{\mathfrak V_s}^2
       =\|f_-\|_{H^s}^2+\|\mathcal Rf_+\|_{H^s}^2
\end{equation}
make $\widehat\Pi_L$ a contraction. The two finite-loss estimates
show that their projective-limit topology is the original smooth
topology. The inverse polarization is not bounded in the unchanged
$L^2$ norms.

For an explicit test, let $L\geq1$ and put
\[
 v_L=\frac{\bar z^L}{(1+|z|^2)^L},\qquad
 \mathcal Cv_L=\frac{z^{L+1}}{(L+1)(1+|z|^2)^{L+1}}.
\]
The second equality follows by differentiation and normalization at
$o$. With $r=|z|^2$, radial integration gives
\[
 \|v_L\|_{L^2}^2
  =\int_0^\infty\frac{r^L}{(1+r)^{2L+2}}\,dr
  =\frac{(L!)^2}{(2L+1)!},\qquad
 \|\mathcal Cv_L\|_{L^2}^2=\frac{(L!)^2}{(2L+3)!}.
\]
Hence
\begin{equation}\label{mht:eq:angular-conditioning}
 \frac{\|\mathcal Cv_L\|_{L^2}^2}{\|v_L\|_{L^2}^2}
       =\frac1{(2L+2)(2L+3)}.
\end{equation}
Since $\mathcal R\mathcal Cv_L=v_L$, these ratios exclude a uniform
$L^2$ bound for $\mathcal R$.

\section{Removal of the cutoff at each fixed arity}

Use the finite complexes
\[
 \widehat E_L^0=\mathbb C1\oplus\widehat V_L,
 \qquad
 \widehat E_L^1=\mathbb C\eta_0\oplus d_E\widehat V_L.
\]
Their inclusion $\widehat i_L$, projection $\widehat p_L$ and
homotopy $\widehat H_L$ are
\begin{align}
 \widehat p_L^0f&=p_0(f)1+
                    \widehat\Pi_L(f-p_0(f)1),\notag\\
 \widehat p_L^1\eta&=\eta_0p_1(\eta)
                        +d_E\widehat\Pi_Lh\eta,&
 \widehat H_L\eta&=(1-\widehat\Pi_L)h\eta.
 \label{mht:eq:angular-contraction}
\end{align}
The homotopy is zero on degree zero. By
Theorem~\ref{mht:thm:finite-cyclic-contraction}, these maps give a
cyclic strong deformation retract preserving the unit and trace.
Their nested homotopies satisfy
\[
 \widehat H_L=\widehat H_M+
                \widehat i_M\widehat H_{ML}\widehat p_M,
 \qquad
 \widehat H_{ML}=(\widehat\Pi_M-\widehat\Pi_L)h|_{\widehat E_M}
 \quad(L\leq M).
\]

Give the coefficient space the graded Sobolev norms
\[
 \mathcal E_s=(E^0\cap H^{s+1})\oplus(E^1\cap H^s),
 \qquad s\geq2.
\]
Theorem~\ref{mht:thm:angular-tame} gives uniform loss six for
$\widehat p_L$ and $\widehat H_L$. Both
$1-\widehat i_L\widehat p_L$ and $\widehat H_L$ have the rapid
tail estimate with this same loss. The additional derivative on
functions matches the gain of $h$ and the derivative in $d_E$.

Let $\beta_{n,L}$ be the shifted operations of
\eqref{mht:eq:finite-tree-recursion} for this contraction.
For a common domain, write $P_L=\widehat i_L\widehat p_L$ and set
\[
 B_{n,L}=\widehat i_L\beta_{n,L}\widehat p_L^{\otimes n}
       :(E[1])^{\otimes n}\longrightarrow E[1].
\]
Suspension is understood on the inclusion and projection. Let
$B_1=b_1$, $B_2=b_2$, and $B_n=0$ for $n\geq3$, using the
signs of Chapter~\ref{mht:ch:finite-reductions}.

\begin{theorem}\label{mht:thm:angular-tree-limit}
For every fixed $n\geq1$, $s\geq2$ and integer $r\geq0$,
\begin{equation}\label{mht:eq:angular-tree-bound}
 \|(B_{n,L}-B_n)(a_1,\ldots,a_n)\|_{\mathcal E_s}
 \leq C_{n,s,r}(1+L)^{-r}
          \prod_{j=1}^n\|a_j\|_{\mathcal E_{s+r+12n+12}}.
\end{equation}
In particular, the finite operations converge uniformly on bounded
smooth sets at each arity. Their binary limit is the native product,
and their higher limits vanish.
\end{theorem}

\begin{proof}
Multiplication is continuous on $\mathcal E_s$. At fixed arity
there are finitely many planar binary trees. When $n\geq3$, each
tree contributing to $\beta_{n,L}$ contains at least one internal
$\widehat H_L$ edge. Use its rapid tail estimate on one edge and
the uniform finite-loss estimates on every other homotopy and
projection. Each root-to-input path meets at most $n-2$ internal
homotopies, one input projection and one output projection.
The displayed derivative budget exceeds their combined losses.

For $n=2$, the extended product is $P_Lb_2(P_L\otimes P_L)$.
Subtract $b_2$ and telescope over the output and the two input
projections. Each difference contains a rapid tail $P_L-1$.
For $n=1$, commutation with $b_1$ gives
$B_{1,L}-b_1=(P_L-1)b_1$, to which the same estimate applies.
This proves \eqref{mht:eq:angular-tree-bound}.

The cyclic and higher-associativity identities involve finitely many
compositions at each arity. The same estimates allow those finite
compositions to pass to the limit.
\end{proof}

The loss grows with arity. Thus all arities together define a formal
series in the field amplitude. The theorem gives no uniform analytic
sum of that series. A fixed finite-order current can be applied to
the smooth output. Products of currents on a common support and
traces around loops require further extension estimates.

\section{Polynomial dependence on the matrix background}

The additional band also removes the artificial kinetic poles of a
fixed compression. Let
\[
 X_L=\widehat\Pi_Lh\eta_0|_{\widehat V_L}.
\]
For $J\in\mathfrak g=\mathfrak{gl}_N(\mathbb C)$, define the
compressed kinetic form by
\[
 K_{J,L}(a,b)=\tau_E\operatorname{tr}
                     \bigl(a(d_Eb+\eta_0[J,b])\bigr),
 \qquad a,b\in\widehat V_L\otimes\mathfrak g.
\]
Write $K_{0,L}$ for the form at $J=0$.

\begin{proposition}\label{mht:prop:angular-background}
The operator $X_L$ raises rotation weight by one and satisfies
$X_L^{2L+2}=0$. Put $N_{J,L}=X_L\otimes\operatorname{ad}_J$.
For every finite $J$,
\begin{align}
 K_{J,L}(a,b)&=K_{0,L}(a,(1+N_{J,L})b),\notag\\
 (1+N_{J,L})^{-1}&=\sum_{q=0}^{2L+1}(-N_{J,L})^q,&
 \det(1+N_{J,L})&=1.
 \label{mht:eq:angular-background-inverse}
\end{align}
Thus the compressed kinetic form is nondegenerate, with inverse
polynomial in the entries of $J$.
\end{proposition}

\begin{proof}
Multiplication by $\eta_0$ raises total rotation weight by one.
Here a one-form includes the weight of $d\bar z$.
The normalized Green operator commutes with rotations fixing $o$,
and the finite corrections defining $h$ occur in the matching weight.
The projection $\widehat\Pi_L$ preserves weights. The weights of
$\widehat V_L$ lie between $-L$ and $L+1$, so $2L+2$ successive
applications of $X_L$ give zero.

For $b\in V$, one has $d_Eh(\eta_0b)=\eta_0b$ because
$p_0(b)=0$. Adjointness of $\widehat\Pi_L$ for $\omega$ gives the
kinetic-form identity. The operator $N_{J,L}$ has the same
nilpotence bound. Its finite geometric series is the inverse, and
all eigenvalues of $1+N_{J,L}$ equal one. Nondegeneracy of
$K_{0,L}$ follows from the cyclic pairing and the full matrix trace.
\end{proof}

For the first four modes, the compressed scalar operator is
\begin{equation}\label{mht:eq:angular-first-background}
 X_1u=-\frac12v,\qquad X_1v=\mathcal Cv,
 \qquad X_1\mathcal Cv=\frac14\mathcal Cu,
 \qquad X_1\mathcal Cu=0.
\end{equation}
Indeed, its matrix is $\Omega^{-1}S$, where $\Omega$ is
\eqref{mht:eq:angular-first-pairing} and
$S_{ab}=\kappa\int v_av_b\,d\mu$. Besides
$\int v^2d\mu=1/12$, the only nonzero entry needed is
$\int u\mathcal Cv\,d\mu=1/24$. These identities give
\eqref{mht:eq:angular-first-background} directly.

The inverse still depends on $J$, despite the unit determinant.
The nilpotence bound alone gives no estimate for that inverse
uniform in $L$. This finite compression therefore leaves the
quantum contraction and cutoff limits as separate questions.

Nor does the rotation grading restore strict associativity. For
$m_2(a,b)=\widehat p_1^0(ab)$ on $\widehat E_1^0$, one has
\[
 m_2(v,v)=\frac1{12},\qquad m_2(v,u)=0,
 \qquad
 m_2(m_2(v,v),u)-m_2(v,m_2(v,u))=\frac u{12}.
\]
The retained weight-zero functions are $1,v$, and the retained
weight-minus-one functions are multiples of $u$.
The moments $\int v^3d\mu=0$ and $\int v|u|^2d\mu=0$ give the
two projected products. The higher operations remain necessary
even for a cutoff whose kinetic form is regular at every background.
